\title{Direct Preference Optimization: Your Language Model is Secretly a Reward Model}

\begin{abstract}
Although large-scale unsupervised language models (LMs) acquire extensive world knowledge and certain reasoning abilities, precisely controlling their behavior remains challenging due to the fully unsupervised nature of their training. Current approaches to achieve such controllability gather human judgments on the relative quality of model outputs and fine-tune the unsupervised LM to align with these preferences, often using reinforcement learning from human feedback (RLHF). Nevertheless, RLHF is a complicated and frequently unstable process, involving initially training a reward model that represents human preferences, and subsequently fine-tuning the large unsupervised LM via reinforcement learning to maximize this predicted reward without deviating excessively from the original model. In this work, we propose a novel parameterization of the reward model in RLHF that allows for the direct extraction of the corresponding optimal policy in closed form, enabling us to address the conventional RLHF task using only a straightforward classification loss. The resulting method, termed \textit{Direct Preference Optimization} (DPO), is stable, efficient, and computationally lightweight, removing the need for sampling from the LM during fine-tuning or extensive hyperparameter tuning. Our empirical results demonstrate that DPO can fine-tune LMs to align with human preferences as well as, or better than, existing techniques. Importantly, DPO-based fine-tuning outperforms PPO-based RLHF in controlling the sentiment of generated outputs and matches or improves the quality of responses in summarization and single-turn dialogue, while being considerably simpler to implement and train.
\end{abstract}

\section{Introduction}  
Large unsupervised language models (LMs) trained on massive datasets develop remarkable abilities~\citep{chowdhery2022palm, brown2020language, touvron2023llama,bubeck2023sparks}. Nevertheless, these models learn from data created by humans with diverse objectives, priorities, and expertise. Some of these human goals and skillsets may not be ideal to replicate; for instance, although we want our AI coding assistant to \textit{recognize} common programming errors to fix them, when generating code, we prefer to guide the model toward the (possibly infrequent) high-quality coding skills present in its training data. Similarly, we might want our language model to be \textit{aware} of a widespread misconception held by 50\% of people, but we certainly do not want the model to assert that misconception as true in half of the related queries! Put differently, carefully choosing the model's \emph{desired responses and behaviors} from its extensive \textit{knowledge and abilities} is essential for creating AI systems that are safe, effective, and controllable \citep{ouyang2022training}. Although existing techniques generally align LMs with human preferences via reinforcement learning (RL), we will demonstrate that the RL-based objective used by these methods can be exactly optimized with a straightforward binary cross-entropy loss, significantly streamlining the preference learning process.

Broadly speaking, current approaches embed desired behaviors into a language model by leveraging curated collections of human preferences that reflect the behaviors humans consider safe and useful. This preference learning phase follows an initial large-scale unsupervised pre-training on extensive text corpora. The simplest way to achieve preference learning is supervised fine-tuning using human demonstrations of high-quality responses. However, the most effective category of methods employs reinforcement learning from human (or AI) feedback (RLHF/RLAIF; \citep{christiano2017deep,bai2022constitutional}). RLHF approaches fit a reward model to datasets of human preference data, then use RL to optimize the language model's policy to generate responses with high reward while avoiding substantial deviation from the original model. Although RLHF yields models with strong conversational and coding performance, the RLHF pipeline is far more complex than supervised learning, involving training multiple language models and repeatedly sampling from the LM policy during training, leading to considerable computational overhead.

In this work, we demonstrate a method to directly optimize a language model to conform to human preferences without relying on explicit reward modeling or reinforcement learning. We introduce \textit{Direct Preference Optimization} (DPO), an algorithm that implicitly targets the same objective as current RLHF algorithms (reward maximization subject to a KL-divergence constraint) while being simple to implement and easy to train. Conceptually, the DPO update boosts the relative log probability of preferred responses over dispreferred ones, but it includes a dynamic, example-specific importance weight that mitigates the model degradation observed with a naive probability ratio objective. Similar to existing approaches, DPO depends on a theoretical preference model (such as the Bradley-Terry model; \cite{bradley1952rankanalysis}) to evaluate how well a given reward function matches empirical preference data. However, whereas current methods employ the preference model to define a preference loss for training a reward model, followed by policy training to optimize this learned reward, DPO applies a change of variables to express the preference loss directly as a function of the policy. Consequently, given a dataset of human preferences over model responses, DPO can optimize a policy via a straightforward binary cross-entropy objective, yielding the optimal policy relative to an implicit reward function fit to the preference data.

Our principal contribution is Direct Preference Optimization (DPO), a straightforward, RL-free algorithm for training language models based on preferences. Our experiments indicate that DPO performs at least as well as established methods, including PPO-based RLHF, for preference learning tasks such as sentiment modulation, summarization, and dialogue, using language models with up to 6B parameters.

\section{Related Work}

Self-supervised language models that scale up in size can perform certain tasks zero-shot \citep{radford2019language} or with few-shot prompts \citep{gpt3,megatron,chowdhery2022palm}. Nonetheless, their effectiveness on downstream tasks and alignment with user intentions can be greatly enhanced by fine-tuning on datasets containing instructions paired with human-generated completions \citep{mishra-etal-2022-cross,sanh2022multitask,chung2022scaling,thoppilan2022lamda}. This process, known as `instruction-tuning,' allows large language models (LLMs) to generalize to instructions beyond those in the fine-tuning dataset, thereby improving their overall usability \citep{chung2022scaling}. Although instruction tuning has proven successful, collecting \textit{relative} human judgments of response quality is often simpler than obtaining expert demonstrations. Consequently, subsequent studies have fine-tuned LLMs on datasets reflecting human preferences, enhancing their capabilities in translation \citep{kreutzer-etal-2018-reliability}, summarization \citep{stiennon2022learning,ziegler2020finetuning}, storytelling \citep{ziegler2020finetuning}, and instruction-following \citep{ouyang2022training,ramamurthy2023is}. These approaches initially train a neural reward model to align with the preference dataset according to a preference model such as the Bradley-Terry model \citep{bradley1952rankanalysis}, and then fine-tune the language model to maximize this reward via reinforcement learning techniques, commonly employing REINFORCE \citep{williams1992reinforce}, proximal policy optimization (PPO; \cite{schulman2017proximal}), or related variants \citep{ramamurthy2023is}. A related research avenue uses LLMs fine-tuned with human feedback for instruction following to produce additional synthetic preference data targeting attributes like safety or harmlessness \citep{bai2022constitutional}, relying solely on weak human supervision through textual rubrics for the LLM’s annotations. These approaches unify two research streams: one focused on training language models using reinforcement learning for various goals~\citep{Ranzato2015SequenceLT,paulus2018a,wu2018learning} and another on general frameworks for learning from human preferences \citep{christiano2017deep,kupcsik2018learning}. Despite the attractiveness of leveraging relative human preferences, fine-tuning large language models via reinforcement learning remains a significant practical hurdle; this work proposes a theoretically grounded method to optimize relative preferences without relying on RL.

Beyond the realm of language, learning policies based on preferences has been explored in both bandit and reinforcement learning frameworks, with a variety of methods proposed. Contextual bandit learning that utilizes preferences or rankings of actions instead of rewards is referred to as a contextual dueling bandit (CDB; \cite{yue2012karmed,dudik2015contextual}). When absolute rewards are unavailable, the theoretical framework for CDBs replaces the concept of an optimal policy with the \textit{von Neumann winner}, defined as a policy whose expected probability of winning against \textit{any} other policy is at least 50\% \citep{dudik2015contextual}. Nevertheless, in the CDB framework, preference labels are obtained online, whereas in learning from human preferences, the learning typically occurs from a fixed offline batch of preference-annotated action pairs \citep{yan2022human}. Likewise, \textit{preference-based RL} (PbRL) uses binary preferences generated by an \textit{unknown} underlying `scoring' function rather than explicit rewards \citep{BusaFekete2014,ruiz2023dueling}. Various PbRL algorithms have been developed, including some that can leverage off-policy preference data, but they generally entail first explicitly estimating the hidden scoring function (i.e., the reward model) and then optimizing it \citep{jain2013learning,BusaFekete2014,christiano2017deep,sadigh2017active,kupcsik2018learning}. In contrast, we propose a single-stage policy learning method that directly optimizes the policy to satisfy preferences.

\section{Preliminaries}\label{section:prelims}

We review the RLHF pipeline as described in \citeauthor{ziegler2020finetuning} (and subsequent works \citep{stiennon2022learning, bai2022training, ouyang2022training}). This process generally consists of three stages: 1) supervised fine-tuning (SFT); 2) preference sampling and reward learning; and 3) reinforcement learning optimization.

\textbf{SFT}: The RLHF process typically starts by fine-tuning a pre-trained language model (LM) via supervised learning on high-quality datasets tailored to the downstream tasks of interest (such as dialogue, summarization, etc.), resulting in a model denoted $\pi^\text{SFT}$. 

\textbf{Reward Modelling Phase}: In the next stage, the SFT model is prompted with inputs $x$ to generate pairs of outputs $(y_1, y_2) \sim \pi^\text{SFT}(y \mid x)$. These pairs are then shown to human annotators, who express a preference for one of the outputs, indicated as $y_w \succ y_l \mid x$, where $y_w$ is the preferred completion and $y_l$ the less preferred among $(y_1, y_2)$. It is assumed that these preferences arise from an underlying latent reward model $r^*(y, x)$, which is not directly observable. Several methods exist to model such preferences, with the Bradley-Terry (BT) \cite{bradley1952rankanalysis} model being a common choice (though broader Plackett-Luce ranking models \citep{plackett1975analysis, luce2012individual} can also be integrated if multiple ranked answers are available). According to the BT model, the human preference distribution $p^*$ can be represented as:

\begin{equation}\label{eq:bradley-terry}
    p^*(y_1\succ y_2 \mid x)=\frac{\exp\left(r^*(x, y_1)\right)}{\exp\left(r^*(x, y_1)\right) + \exp\left(r^*(x, y_2)\right)}.
\end{equation}

Given access to a fixed dataset of comparisons $\mathcal{D}=\bigl\{x^{(i)}, y_w^{(i)}, y_l^{(i)}\bigr\}_{i=1}^N$ drawn from $p^*$, we can model a reward function $r_{\phi}(x, y)$ and estimate its parameters through maximum likelihood estimation. By casting the problem as a binary classification task, we obtain the negative log-likelihood loss:

\begin{equation}\label{eq:reward_model}
    \mathcal{L}_R(r_{\phi}, \mathcal{D}) = -\mathbb{E}_{(x, y_w, y_l)\sim \mathcal{D}}\bigl[\log \sigma(r_{\phi}(x, y_w)- r_{\phi}(x, y_l))\bigr]
\end{equation}

where $\sigma$ denotes the logistic function. In the setting of language models (LMs), the reward network $r_{\phi}(x, y)$ is frequently initialized using the SFT model $\pi^\text{SFT}(y \mid x)$, augmented with a linear layer atop the final transformer layer that outputs a scalar reward prediction \cite{ziegler2020finetuning}. To reduce variance in the reward values, previous work normalizes the rewards such that $\mathbb{E}_{x,y\sim \mathcal{D}}\left[r_\phi(x, y)\right] = 0$ for every $x$.

\textbf{Reinforcement Learning Fine-Tuning Phase}: During the RL fine-tuning stage, the learned reward function is employed to guide the language model’s updates. Following earlier studies~\citep{jaques2017sequence, jaques2020human}, the optimization objective is expressed as

\begin{equation}\label{eq:RL}
\max_{\pi_{\theta}}  \mathbb{E}_{x\sim \mathcal{D}, y\sim \pi_{\theta}(y \mid x)}\bigl[r_{\phi}(x, y)\bigr] - \beta\mathbb{D}_{\textrm{KL}}\bigl[\pi_{\theta}(y\mid x)\mid \mid \pi_\text{ref}(y\mid x)\bigr],
\end{equation}

where $\beta$ regulates the degree to which the updated policy $\pi_\theta$ deviates from the baseline reference policy $\pi_\text{ref}$, typically the initial SFT model $\pi^\text{SFT}$. In practice, the policy $\pi_\theta$ is initialized to $\pi^\text{SFT}$ as well. This constraint is crucial to prevent the policy from straying far from the distribution where the reward model’s predictions are reliable, while also preserving diversity in generation and avoiding collapse to a narrow set of high-reward outputs. Because language generation is discrete, this objective is not differentiable and is generally optimized using reinforcement learning algorithms. The common strategy \citep{ziegler2020finetuning, stiennon2022learning, bai2022training, ouyang2022training} involves defining the reward function as ${r(x, y) = r_{\phi}(x, y) -\beta (\log \pi_{\theta}(y\mid x) - \log \pi_\text{ref}(y\mid x))}$, then maximizing it using PPO \cite{schulman2017proximal}.

\section{Direct Preference Optimization}\label{sec:DPO}

To address the difficulties in applying reinforcement learning techniques to large-scale tasks such as language model fine-tuning, we aim to develop a straightforward method for policy optimization directly from preferences. Unlike previous RLHF approaches, which first learn a reward function and then optimize it via reinforcement learning, our method adopts a specific parameterization of the reward model that allows for derivation of its optimal policy in closed form, obviating the need for an RL training loop. As we will elaborate, the core insight is to exploit an analytical relationship mapping reward functions to their optimal policies, thereby transforming a loss defined over reward functions into one defined over policies. This change-of-variable technique eliminates the requirement to explicitly fit a standalone reward model while still optimizing according to established human preference models such as the Bradley-Terry model. Effectively, the policy network simultaneously

\textbf{Deriving the DPO objective.} We begin with the same reinforcement learning (RL) objective as previous studies, Eq.~\ref{eq:RL}, using a general reward function $r$. Following earlier work~\citep{peters2007reinforcement, peng2019advantage, korbak2022reinforcement, go2023aligning}, it is straightforward to demonstrate that the optimal solution to the KL-constrained reward maximization problem in Eq.~\ref{eq:RL} is given by:

\begin{equation}\label{eq:op_policy}
    \pi_r(y\mid x) = \frac{1}{Z(x)}\pi_\text{ref}(y\mid x)\exp\left(\frac{1}{\beta}r(x, y)\right),
\end{equation}

where the partition function is defined as $Z(x) =\sum_{y}\pi_\text{ref}(y\mid x)\exp\left(\frac{1}{\beta}r(x, y)\right)$. A full derivation can be found in Appendix \ref{app:derivation1}. Even when employing the MLE estimate $r_{\phi}$ for the true reward function $r^*$, estimating the partition function $Z(x)$ remains computationally demanding \citep{korbak2022reinforcement, go2023aligning}, rendering this form challenging to apply practically. Nonetheless, we can rearrange Eq.~\ref{eq:op_policy} to rewrite the reward function in terms of its corresponding optimal policy $\pi_r$, the reference policy $\pi_\text{ref}$, and the unknown partition function $Z(\cdot)$. Specifically, by taking the logarithm of both sides of Eq.~\ref{eq:op_policy} and performing some algebraic manipulation, we get:

\begin{equation}\label{eq:main_eq}
    r(x,y) =\beta \log \frac{\pi_r(y\mid x)}{\pi_\text{ref}(y\mid x)} + \beta \log Z(x).
\end{equation}

This reparameterization can be applied to the true reward $r^*$ and its associated optimal policy $\pi^*$. Fortunately, the Bradley-Terry model only depends on the difference in rewards between two completions, i.e., ${p^*(y_1 \succ y_2 \mid x) = \sigma(r^*(x, y_1) - r^*(x, y_2))}$. By substituting the expression from Eq.~\ref{eq:main_eq} for $r^*(x,y)$ into the preference model Eq.~\ref{eq:bradley-terry}, the partition function terms cancel out, allowing us to express the human preference probability solely in terms of the optimal policy $\pi^*$ and the reference policy $\pi_\text{ref}$. Hence, the optimal RLHF policy $\pi^*$ under the Bradley-Terry framework satisfies the preference model:

\begin{equation}\label{eq:objective}
    p^*(y_1\succ y_2 \mid x)=\frac{1}{1 + \exp\left(\beta \log \frac{\pi^*(y_2\mid x)}{\pi_\text{ref}(y_2\mid x)} - \beta \log \frac{\pi^*(y_1\mid x)}{\pi_\text{ref}(y_1\mid x)}\right)}.
\end{equation}

The full derivation is provided in Appendix~\ref{app:derivation2}. Although Eq.~\ref{eq:objective} is derived using the Bradley-Terry model, similar derivations can be made for more general Plackett-Luce models~\citep{plackett1975analysis, luce2012individual}, as detailed in Appendix~\ref{app:plackett_luce_models}.

Having expressed the probability of human preference data in terms of the optimal policy rather than the reward model, we can now define a maximum likelihood objective for a parametrized policy $\pi_\theta$. Similar to the reward modeling objective (Eq.~\ref{eq:reward_model}), our policy objective becomes:

\begin{equation}\label{eq:optimum_model}
    \mathcal{L}_\text{DPO}(\pi_{\theta}; \pi_\text{ref}) = -\mathbb{E}_{(x, y_w, y_l)\sim \mathcal{D}}\left[\log \sigma \left(\beta \log \frac{\pi_{\theta}(y_w\mid x)}{\pi_\text{ref}(y_w\mid x)} - \beta \log \frac

In this manner, we fit an implicit reward through an alternative parameterization, whose optimal policy corresponds directly to $\pi_\theta$. Additionally, since our method is equivalent to fitting a reparametrized Bradley-Terry model, it inherits certain theoretical advantages, including consistency under appropriate assumptions about the preference data distribution \cite{bong2022generalized}. In Section~\ref{sec:theory}, we further elaborate on the theoretical aspects of DPO in comparison with related studies.

\textbf{What effect does the DPO update have?} To gain a mechanistic insight into DPO, it is helpful to examine the gradient of the loss function $\mathcal{L}_\text{DPO}$. The gradient with respect to the parameters $\theta$ can be expressed as:

\begin{multline*}\label{eq:gradient}
    \nabla_\theta \mathcal{L}_\text{DPO}(\pi_\theta;\pi_\text{ref}) = \\ -\beta\mathbb{E}_{(x, y_w, y_l) \sim \mathcal{D}} \bigg[\underbrace{\sigma(\hat{r}_\theta(x, y_l) - \hat{r}_\theta (x, y_w))}_\text{higher weight when reward estimate is wrong}\bigg[\underbrace{\nabla_\theta\log \pi(y_w \mid x)}_\text{increase likelihood of $y_w$} - \underbrace{\nabla_\theta\log\pi(y_l \mid x)}_\text{decrease likelihood of $y_l$}\bigg]\bigg],
\end{multline*}

where $\hat{r}_\theta(x, y) = \beta \log \frac{\pi_\theta(y \mid x)}{\pi_\text{ref}(y \mid x)}$ denotes the reward implicitly defined by the language model $\pi_\theta$ and the reference model $\pi_\text{ref}$ (further details in Section~\ref{sec:theory}). Conceptually, the gradient of $\mathcal{L}_\text{DPO}$ serves to raise the probability of preferred completions $y_w$ while lowering that of the less preferred completions $y_l$. Crucially, the updates are weighted by the extent to which the implicit reward model $\hat{r}_\theta$ assigns a higher score to the disfavored completions, scaled by $\beta$—that is, how inaccurately the implicit reward model ranks these completions, taking into account the strength of the KL constraint. Our empirical findings underscore the significance of this weighting; a naive variant of this approach lacking the weighting factor can lead to language model degeneration (see Appendix Table~\ref{tab:unlikelihood_generations}).

\textbf{DPO overview.} The standard DPO workflow proceeds as follows: 1) For each prompt $x$, generate completions $y_1, y_2 \sim \pi_\text{ref}(\cdot \mid x)$, then label these using human preferences to form the offline preference dataset $\mathcal{D} = \{x^{(i)}, y_w^{(i)}, y_l^{(i)}\}_{i=1}^N$; and 2) train the language model $\pi_\theta$ by minimizing the loss $\mathcal{L}_\text{DPO}$ given the fixed $\pi_\text{ref}$, dataset $\mathcal{D}$, and chosen $\beta$. In practical scenarios, it is preferable to utilize publicly available preference datasets rather than generating new samples and collecting fresh human judgments. Since these preference datasets are typically sampled using $\pi^\text{SFT}$, we set $\pi_\text{ref} = \pi^\text{SFT}$ when possible. If $\pi^\text{SFT}$ is unavailable, we instead initialize $\pi_\text{ref}$ by maximizing the likelihood of the preferred completions $(x, y_w)$, specifically, ${\pi_\text{ref} = \argmax_{\pi} \mathbb{E}_{x, y_w \sim \mathcal{D}}\left[\log \pi(y_w \mid x)\right]}$. This approach helps reduce the distribution shift between the inaccessible true reference distribution and the $\pi_\text{ref}$ employed in DPO. Additional implementation details and hyperparameter settings are provided in Appendix~\ref{app:implementation}.

\section{Theoretical Analysis of DPO} This section offers a deeper interpretation of the DPO approach, furnishes theoretical justification, and connects the benefits of DPO to limitations encountered by actor-critic methods used in RLHF (e.g., PPO~\cite{schulman2017proximal}).

\label{sec:theory}

\subsection{Your Language Model Is Secretly a Reward Model} DPO circumvents the need to explicitly fit a reward function or carry out reinforcement learning by optimizing a single maximum likelihood objective. Observe that the objective in Eq. \ref{eq:main_eq} corresponds to a Bradley-Terry model where the reward function is parameterized as $r^*(x, y) = \beta \log \frac{\pi^*_\theta(y \mid x)}{\pi_\text{ref}(y \mid x)}$. Training our parametric model $\pi_\theta$ is therefore equivalent to optimizing a reward model as described in Eq. \ref{eq:reward_model} after an appropriate change of variables. In this subsection, we develop the theoretical foundations underpinning this reparameterization, demonstrate that it does not restrict the set of learnable reward models, and show it enables exact recovery of the optimal policy. We begin by introducing an equivalence relation among reward functions.

\begin{definition}
We define two reward functions $r(x, y)$ and $r'(x, y)$ to be equivalent if and only if there exists some function $f$ such that ${r(x, y)-r'(x, y) = f(x)}$.     
\end{definition}

It is straightforward to verify that this relation is indeed an equivalence relation, which divides the set of reward functions into distinct classes. We can then formulate the following two lemmas:

\begin{lemma}\label{lemma:same_prefrence} Within the Plackett-Luce framework, including specifically the Bradley-Terry model, any two reward functions belonging to the same equivalence class generate identical preference distributions.
\end{lemma}

\begin{lemma}\label{lemma:same_policy}
    Two reward functions in the same equivalence class lead to the same optimal policy for the constrained reinforcement learning problem.
\end{lemma}

The proofs are direct and are postponed to Appendix \ref{app:lemma1}. The first lemma reflects a well-known under-specification characteristic of the Plackett-Luce family \cite{plackett1975analysis}. Because of this under-specification, it is common to impose extra identifiability constraints to guarantee the properties of MLE estimates from Eq. \ref{eq:reward_model} \cite{bong2022generalized}. The second lemma implies that all reward functions within one equivalence class produce the same optimal policy, so for our ultimate goal, it suffices to recover any reward function from the relevant class. The following theorem is proved in Appendix~\ref{app:thm1}:

\begin{theorem}\label{thm:main}
    Under mild conditions, every reward class compatible with Plackett-Luce models (and Bradley-Terry in particular) can be expressed in the form ${r(x, y) = \beta \log \frac{\pi(y\mid x)}{\pi_\text{ref}(y\mid x)}}$, for some model $\pi(y\mid x)$ and a given reference model $\pi_\text{ref}(y \mid x)$.
\end{theorem}

\begin{sproof}
    Take any reward function $r(x, y)$, which corresponds to an optimal model $\pi_r(y \mid x)$ defined by Eq. \ref{eq:op_policy}. We will demonstrate that a reward function equivalent to $r$ can be written using the above reparameterization. Define the projection $f$ as  
\begin{equation}
    f(r; \pi_\text{ref}, \beta)(x, y) = r(x, y) - \beta\log\sum_{y}\pi_\text{ref}(y\mid x)\exp\left(\frac{1}{\beta}r(x, y)\right)
\end{equation}
The operator $f$ effectively normalizes the reward function by subtracting the logarithm of the partition function of $\pi_r$. Since this normalization term depends only on the prefix $x$, $f(r; \pi_\text{ref}, \beta)(x, y)$ lies in the equivalence class of $r(x, y)$. Finally, substituting $r$ with the expression on the right side of Eq.~\ref{eq:main_eq} (which holds for any reward function), we get $f(r; \pi_\text{ref}, \beta)(x, y) = \beta \log \frac{\pi_r(y\mid x)}{\pi_\text{ref}(y\mid x)}$. Hence, the projection $f$ yields a member of $r$'s equivalence class in the desired form, ensuring that the proposed reparameterization does not restrict the generality of our reward model.
\end{sproof}

An alternative interpretation of Theorem~\ref{thm:main} is that it precisely identifies which reward function within each equivalence class is chosen by the DPO reparameterization, specifically the reward function that satisfies:

\begin{equation}\label{eq:lag_p}
     \sum_{y}\underbrace{\pi_\text{ref}(y\mid x)\exp\left(\frac{1}{\beta}r(x, y)\right)}_{=\pi(y\mid x)\text{, using Thm.~\ref{thm:main} reparam.}} = 1,
\end{equation}

meaning that $\pi(y\mid x)$ constitutes a valid probability distribution (with positive values summing to one). Moreover, as indicated by Eq.~\ref{eq:op_policy}, Eq.~\ref{eq:lag_p} corresponds to the partition function of the optimal policy derived from the reward function $r(x, y)$. The principal insight behind the DPO method is that it enforces certain constraints on the underdetermined Plackett-Luce family (especially Bradley-Terry models) of preference models, preserving the class of reward functions that can be represented, while ensuring that the optimal policy expressed in Eq.~\ref{eq:op_policy} is analytically manageable for every prompt $x$.

\subsection{Instability of Actor-Critic Algorithms} Our framework can also be employed to analyze instabilities observed in conventional actor-critic algorithms commonly used in RLHF, such as PPO. By following the RLHF procedure and concentrating on the RL fine-tuning phase described in Section \ref{section:prelims}, we establish links to the control-as-inference framework \cite{levine2018reinforcement} for the constrained RL problem presented in \ref{eq:RL}. We consider a parameterized policy model $\pi_{\theta}(y\mid x)$ and aim to minimize the divergence $\mathbb{D}_{\text{KL}}[\pi_{\theta}(y|x) \mid \mid \pi^*(y\mid x)]$, where $\pi^*$ denotes the optimal policy given by Eq.~\ref{eq:optimum_model}, derived from the reward function $r_{\phi}(y, x)$. Through algebraic manipulation, this yields the following optimization objective:

\begin{equation}\label{eq:AC}
    \max_{\pi_{\theta}}\mathbb{E}_{\pi_{\theta}(y\mid x)}\bigg[\underbrace{r_{\phi}(x, y) -\beta\log\sum_{y}\pi_\text{ref}(y\mid x)\exp\left(\frac{1}{\beta}r_{\phi}(x, y)\right)}_{f(r_{\phi}, \pi_\text{ref}, \beta)} - \underbrace{\beta\log\frac{\pi_{\theta}(y\mid x)}{\pi_\text{ref}(y\mid x)}}_{\text{KL}}\bigg]
\end{equation}

This objective is identical to the one optimized in previous studies \citep{ziegler2020finetuning, stiennon2022learning, bai2022training, ouyang2022training} using the DPO-equivalent reward for the reward class $r_{\phi}$. Within this framework, the normalization term in $f(r_{\phi}, \pi_\text{ref}, \beta)$ can be viewed as the soft value function of the reference policy $\pi_\text{ref}$. Although this term does not influence the optimal solution, omitting it can result in a policy gradient with high variance, which may cause instability in learning. One approach to handle the normalization term is to use a learned value function, but optimizing that can be challenging. Alternatively, previous methods have normalized rewards by employing a human completion baseline, effectively a single-sample Monte Carlo estimate of the normalization term. In contrast, the DPO reparameterization produces a reward function that does not require any baseline adjustments.

\section{Experiments}
This section presents an empirical evaluation of DPO's capability to train policies directly from preference data. Initially, in a controlled text-generation environment, we investigate how effectively DPO balances maximizing reward with minimizing KL-divergence relative to the reference policy, compared to widely used preference learning methods such as PPO. Subsequently, we test DPO on larger-scale models and more complex RLHF tasks, including summarization and dialogue. We observe that, with minimal hyperparameter tuning, DPO generally matches or surpasses the performance of robust baselines like RLHF with PPO and also outperforms methods that return the best trajectory among $N$ samples under a learned reward function. Prior to detailing these findings, we outline the experimental setup; further specifics are provided in Appendix~\ref{app:exp_details}.

\textbf{Tasks.} Our study investigates three distinct open-ended text generation tasks. In all cases, the algorithms acquire a policy based on a preference dataset $\mathcal{D}=\bigl\{x^{(i)}, y_w^{(i)}, y_l^{(i)}\bigr\}_{i=1}^N$. In \textbf{controlled sentiment generation}, the input $x$ represents the beginning segment of a movie review from the IMDb dataset \cite{maas-EtAl:2011:ACL-HLT2011}, and the objective of the policy is to generate $y$ exhibiting positive sentiment. To enable controlled evaluation, for this task we \textit{create} preference pairs from generated outputs using a pre-trained sentiment classifier, ensuring that $p(\text{positive}\mid x,y_w) > p(\text{positive}\mid x,y_l)$. For SFT, we fine-tune GPT-2-large on the training portion of the IMDb reviews until convergence (additional information in App~\ref{app:sentiment_details}). In \textbf{summarization}, the input $x$ is a forum post from Reddit; the policy must produce a summary $y$ capturing the key points of the post. Consistent with prior studies, we utilize the Reddit TL;DR summarization dataset \citep{volske-etal-2017-tl} along with human preference data collected by \citeauthor{stiennon2022learning}. We employ an SFT model fine-tuned on human-generated forum post summaries\footnote{\url{https://huggingface.co/CarperAI/openai_summarize_tldr_sft}} using the TRLX \citep{leandro_von_werra_2023_7790115} framework for RLHF. The human preference dataset was sourced by \citeauthor{stiennon2022learning} from samples of a distinct but similarly fine-tuned SFT model. Lastly, in \textbf{single-turn dialogue}, the input $x$ corresponds to a human query, which could range from astrophysics questions to requests for relationship advice. The policy is tasked with generating a response $y$ that is both engaging and helpful to the user; here we utilize the Anthropic Helpful and Harmless dialogue dataset \citep{bai2022training}, which includes 170k conversations between humans and an automated assistant. Each dialogue concludes with a pair of responses generated by a large (though undisclosed) language model, along with a preference label indicating the human-favored reply. In this context, as no pre-trained SFT model exists, we fine-tune a standard language model exclusively on the preferred completions to create the SFT model.

\textbf{Evaluation.} Our experiments employ two distinct evaluation methodologies. To assess how effectively each algorithm optimizes the constrained reward maximization objective within the controlled sentiment generation context, we measure each algorithm's performance by examining the frontier formed by the reward achieved and the KL-divergence from the reference policy; this frontier can be computed since the true reward function (a sentiment classifier) is accessible. Conversely, in practical scenarios where the true reward function is unknown, we assess algorithms based on their \textit{win rate} against a baseline policy, utilizing GPT-4 as a stand-in for human evaluation of summary quality and response helpfulness in the summarization and single-turn dialogue tasks, respectively. For summarization, we adopt the reference summaries in the test set as the baseline; for dialogue, the preferred response in the test dataset serves as the baseline. Although prior research indicates that language models can serve as superior automated evaluators compared to traditional metrics \citep{Chen2023ExploringTU}, we perform a human study to validate our use of GPT-4 for evaluation, detailed in Sec.~\ref{sec:human-judgments}. Our findings show that GPT-4's judgments correlate strongly with human assessments, with agreement between humans and GPT-4 typically matching or exceeding inter-human annotator agreement. 

\textbf{Methods.} Alongside DPO, we assess multiple existing strategies for training language models to align with human preferences. As a straightforward approach, we explore zero-shot prompting with \textbf{GPT-J} \citep{gpt-j} for the summarization task and 2-shot prompting with \textbf{Pythia-2.8B} \citep{biderman2023pythia} for the dialogue task. Additionally, we evaluate the \textbf{SFT} model and \textbf{Preferred-FT}, which is fine-tuned via supervised learning on the chosen completion $y_w$ originating either from the SFT model (in controlled sentiment and summarization) or from a general LM (in single-turn dialogue). Another pseudo-supervised technique is \textbf{Unlikelihood}~\citep{welleck2019neural}, which trains the policy to increase the probability assigned to $y_w$ while simultaneously \textit{decreasing} the probability allocated to $y_l$; this method employs an optional coefficient $\alpha \in [0,1]$ on the unlikelihood component. We also include \textbf{PPO} \citep{schulman2017proximal} trained using a reward function derived from the preference data, and \textbf{PPO-GT}, an oracle variant that learns from the true reward function available in the controlled sentiment setting. For sentiment tasks, we implement two versions of PPO-GT: an off-the-shelf version \cite{leandro_von_werra_2023_7790115} and a modified one that normalizes rewards and further tunes hyperparameters for enhanced performance (these adjustments are similarly applied when executing standard PPO with learned rewards). Lastly, we consider the \textbf{Best of $N$} baseline, which involves sampling $N$ responses from the SFT model (or Preferred-FT for dialogue) and selecting the highest-scoring response based on a reward model trained on the preference dataset. Although this high-performing technique separates the reward model quality from the PPO optimization process, it is computationally intensive for even moderate values of $N$ because it requires generating $N$ completions for each query at test time.

\subsection{How effectively can DPO optimize the RLHF objective?}

The KL-constrained reward maximization objective commonly employed in RLHF algorithms balances maximizing reward with limiting the policy's divergence from the reference policy. Thus, when comparing different algorithms, it is essential to consider both the achieved reward and the KL divergence; obtaining a slightly higher reward at the expense of a significantly larger KL divergence is not necessarily preferable. Figure~\ref{fig:frontier-tldr-main} illustrates the reward-KL frontier for various algorithms in the sentiment task. We conduct multiple training runs for each algorithm, varying the hyperparameter controlling policy conservativeness in each run (target KL $\in\{3,6,9,12\}$ for PPO, $\beta \in \{0.05,0.1,1,5\}$, $\alpha\in\{0.05,0.1,0.5,1\}$ for unlikelihood, and different random seeds for preferred-FT). This sweep comprises a total of 22 runs. After every 100 training steps until convergence, we evaluate each policy on a set of test prompts, calculating the average reward according to the true reward function as well as the average sequence-level KL\footnote{Namely, the sum of the KL divergences at each timestep.} relative to the reference policy, $\text{KL}\left(\pi\mid \mid \pi_\text{ref}\right)$. Our findings show that DPO yields the most efficient frontier by a significant margin, attaining the highest reward while maintaining a low KL divergence. This outcome is particularly noteworthy for several reasons. First, although DPO and PPO optimize the same objective, DPO is considerably more efficient; the tradeoff between reward and KL for DPO strictly outperforms that of PPO. Second, DPO surpasses PPO's frontier even \emph{when PPO has access to ground truth rewards} (PPO-GT).

\subsection{Is DPO scalable to real-world preference datasets?}
\label{sec:dpo-real-datasets}
Next, we assess the fine-tuning effectiveness of DPO on tasks involving summarization and single-turn dialogue. In the context of summarization, automatic evaluation metrics like ROUGE often show weak correlation with human preferences~\citep{stiennon2022learning}, and prior research has demonstrated that fine-tuning language models (LMs) with PPO based on human preferences can yield better summaries. We compare various methods by generating completions from the test split of the TL;DR summarization dataset and calculating the average win rate against the reference completions in the test set. For all methods, completions are sampled with temperatures ranging from 0.0 to 1.0, and the resulting win rates are depicted in Figure~\ref{fig:frontier-tldr-main} (right). DPO, PPO, and Preferred-FT all fine-tune the same GPT-J SFT model\footnote{\url{https://huggingface.co/CarperAI/openai_summarize_tldr_sft}}. We observe that DPO achieves a win rate of about 61\% at temperature 0.0, surpassing PPO’s performance of roughly 57\% at its optimal temperature of 0.0. Furthermore, DPO attains a higher peak win rate than the best of $N$ baseline. It is important to mention that we did not extensively tune DPO’s $\beta$ hyperparameter, so these results might underestimate DPO’s capabilities. Additionally, DPO demonstrates much greater robustness to variations in sampling temperature compared to PPO, whose performance can deteriorate to that of the original GPT-J model at higher temperatures. Preferred-FT shows minimal improvement over the SFT model. We also present a direct comparison of DPO and PPO through human evaluations in Section~\ref{sec:human-judgments}, where samples generated by DPO at temperature 0.25 were preferred 58\% of the time over those produced by PPO at temperature 0.

For the single-turn dialogue task, we assess various methods using the subset of the Anthropic HH dataset test split \citep{bai2022training} that involves one round of human-assistant interaction. GPT-4 evaluations utilize the preferred completions from the test set as a benchmark to calculate the win rate across different methods. Since there is no established SFT model for this task, we begin with a pre-trained Pythia-2.8B model, apply Preferred-FT to train a reference model on the selected completions—ensuring these completions remain within the model's distribution—and subsequently train with DPO. We also benchmark against the best among 128 Preferred-FT completions (noting that the Best of $N$ baseline saturates at 128 completions for this task; see Appendix Figure~\ref{fig:best-of-n}) and a 2-shot prompted version of the Pythia-2.8B base model, finding that DPO matches or exceeds performance at the optimal temperatures for each method. Additionally, we evaluate an RLHF model trained with PPO on the Anthropic HH dataset\footnote{\url{https://huggingface.co/reciprocate/ppo_hh_pythia-6B}} from a recognized source\footnote{\url{https://github.com/CarperAI/trlx/tree/main/examples/hh}}, but fail to identify any prompt or sampling temperature that outperforms the base Pythia-2.8B model. Drawing from our TL;DR results and the fact that both approaches optimize the same reward function, we regard Best of 128 as an approximate indicator of PPO-level performance. In summary, DPO stands out as the only computationally efficient approach that surpasses the preferred completions on the Anthropic HH dataset and achieves comparable or superior results relative to the computationally intensive Best of 128 baseline. Lastly, Figure~\ref{fig:dialogue-main} illustrates that DPO attains its peak performance relatively swiftly.

\subsection{Generalization to a new input distribution}

To further assess the performance differences between PPO and DPO under distributional shifts, we tested the PPO and DPO policies derived from our Reddit TL;DR summarization experiment on a different dataset—news articles from the test split of the CNN/DailyMail dataset \citep{nallapati-etal-2016-abstractive}. We used the optimal sampling temperatures identified on TL;DR (0 and 0.25). The outcomes are shown in Table~\ref{tab:ood}. We calculated the GPT-4 win rate against the ground-truth summaries in the datasets using the same GPT-4 (C) prompt employed for Reddit TL;DR, but substituting the phrase ``forum post'' with ``news article''. On this new distribution, DPO still significantly outperforms the PPO policy. This experiment offers preliminary evidence that DPO policies can generalize as effectively as PPO policies, despite DPO not leveraging the additional unlabeled Reddit TL;DR prompts used by PPO.

\subsection{Validating GPT-4 judgments with human judgments}
\label{sec:human-judgments}

We carried out a human evaluation study to verify the dependability of GPT-4’s assessments, utilizing the results from the TL;DR summarization task and two distinct GPT-4 prompts. The \textbf{GPT-4 (S)} (simple) prompt asks directly which summary better captures the key information of the post. The \textbf{GPT-4 (C)} (concise) prompt additionally inquires which summary is more concise; we evaluate this prompt because we observed that GPT-4 tends to favor longer, more repetitive summaries with the \textbf{GPT-4 (S)} prompt than humans do. The full prompts are provided in Appendix~\ref{app:prompts}. We perform three pairwise comparisons using methods that represent high (DPO, temp. 0.25), low (PPO, temp. 1.0), and intermediate (SFT, temp. 0.25) performance levels, aiming to cover a range of sample qualities; all three are compared to greedily sampled PPO (its best temperature). We find that for both prompts, GPT-4 agrees with human raters about as frequently as humans agree with one another, indicating that GPT-4 serves as a reasonable stand-in for human evaluations (due to a limited number of human raters, multiple human judgments were only collected for the DPO and PPO-1 comparisons). Overall, the \textbf{GPT-4 (C)} prompt yields win rates that better reflect human preferences; thus, we use this prompt for the primary results reported in Section~\ref{sec:dpo-real-datasets}. Further details on the human study, including the interface shown to raters and a list of participants, can be found in Appendix~\ref{app:human-study}.

\section{Discussion}
Learning from preferences offers a robust and scalable approach to developing capable, aligned language models. We have presented DPO, a straightforward training method that enables language models to learn from preferences without relying on reinforcement learning. Instead of forcing the preference learning challenge into a conventional RL framework to apply standard RL algorithms, DPO establishes a correspondence between language model policies and reward functions, allowing the language model to be trained to directly align with human preferences using a simple cross-entropy loss—eliminating the need for reinforcement learning and avoiding any loss of generality. With minimal hyperparameter tuning, DPO achieves comparable or superior performance relative to existing RLHF methods, including PPO-based approaches, significantly lowering the barrier for training language models from human preferences.

\textbf{Limitations \& Future Work.} Our findings prompt several key questions for subsequent research. How does the DPO policy generalize beyond the training distribution in comparison to models trained via explicit reward functions? Preliminary evidence indicates that DPO policies generalize similarly to those trained with PPO, but more extensive investigation is required. For instance, could self-labeling with the DPO policy also enable efficient utilization of unlabeled prompts? Moreover, how does reward over-optimization appear within the direct preference optimization framework, and could the modest performance drop seen in Figure~\ref{fig:dialogue-main}-right be an example of this? Although we evaluate models with up to 6B parameters, extending DPO to scale to cutting-edge models with significantly larger parameter counts represents a promising avenue for future exploration. Regarding evaluation metrics, we observe that GPT-4’s win rates are influenced by the prompt design; further research could determine optimal methods for eliciting reliable judgments from automated evaluators. Lastly, DPO has a broad range of potential applications beyond training language models from human preferences, such as training generative models in various other modalities.

\end{document}